% NAF 4073, batch 2 — Gallica views 21–40.
% Pass: Opus 5 (claude-opus-5), 2026-09-11 — first pass, unchecked against
% the views by a human.
%
% What the twenty views hold. The batch carries folios 13v to 31r of the
% volume: eleven pieces in six hands. Three more Fourier billets close the
% run begun in batch 1; then, at ff. 20–21, the one piece of this batch in
% Sophie Germain's own hand — a long letter to Libri dated « Paris ce 18
% avril 1831 », two months before her death, which speaks of her illness, of
% Fourier's death and of « cet élève Gallois ». After it come a note of
% Lagrange in the third person, three letters of Lagrange written from
% Berlin to a correspondent at Turin (not to Germain), a note of Lalande of
% 4 nov. 1797, a letter of Libri from Florence of 17 November 1825, and the
% opening leaf of an unsigned letter that runs on past the batch.
%
%   v. 21     f. 13v?   The address leaf of the Fourier billet « mardi soir »
%                       transcribed at view 19 in batch 1. No pencilled folio
%                       is legible on this leaf.
%   v. 21     f. 14r    Fourier, « jeudi … mai » (a word struck between the
%                       two): he invites Germain and her mother to lunch on
%                       the Sunday with the Pronys, « à midi un quart », and
%                       sends an engraved portrait of « un de nos
%                       prédécesseurs ».
%   v. 23     f. 15v    The address leaf of that letter.
%   v. 23     f. 16r    Fourier, « mardi matin »: he will come to the dinner
%                       of the … December, and « il me tarde beaucoup
%                       mademoiselle de reprendre nos entretiens
%                       analytiques ». He apologises for praising her,
%                       « parce que les éloges vous déplaisent ».
%   v. 25     f. 17v    The address leaf of that letter.
%   v. 25     f. 18r    Fourier, « lundi matin »: he forwards a note from
%                       M. de Prony and will fetch from the Institut the
%                       copy of Libri's memoir he could not find at home.
%   v. 27     f. 19v    The address leaf of that letter.
%   v. 27–29  ff. 20r–  SOPHIE GERMAIN to Libri, « Paris ce 18 avril 1831 »,
%             21v       signed « Sophie Germain » — four written pages, the
%                       only leaves of this batch in her hand. Her health
%                       (« une mort prompte seroit pour moi un soulagement »),
%                       the parcels that did not arrive, Mme Cauchy, Mme
%                       Renaud, her sister; then, in one paragraph, the state
%                       of mathematics — « votre préoccupation celle de
%                       Cauchy, la mort de Mr fourier » — and Galois, whom
%                       she spells GALLOIS: expelled from the École normale,
%                       his mother reduced to placing herself as a lady's
%                       companion, « on dit qu'il devient tout à fait fou et
%                       je le crains ».
%   v. 29     f. 22r    Lagrange, in the third person, « Paris ce 17 germinal
%                       (mercredi) »: he will be at Mlle Germaine's orders on
%                       the 19th and the 20th. He writes her name GERMAINE.
%   v. 31     f. 23v    The address leaf of that note.
%   v. 31–32  ff. 24r–  Lagrange, « à Berlin ce 10 Juillet 1776 », signed
%             24v       « De la Grange » — NOT to Germain: to a correspondent
%                       at Turin, of Denina, of « le Docteur Cigna », of
%                       « Ms. Richeri, secrétaire de S. E. le Comte
%                       Fontana ». The opening leaf is wanting: the text
%                       begins « Voici deux lettres… » with no salutation.
%   v. 32–33  ff. 25r–  Lagrange, same hand, neither salutation nor
%             25v       signature: on Denina's troubles, on conforming to the
%                       laws of the country one lives in, and on Galileo,
%                       whose Dialogues « sont les moins bons de tous ses
%                       ouvrages ». Ends by sending « une medaille de
%                       l'Imperatrice de Russie ».
%   v. 33–34  ff. 26r–  Lagrange, same hand, opening « Monsieur », no
%             26v       signature: the engravings of the King sent by
%                       M. Verney, Weguelin's diplomatic history, and the
%                       armies in the field — « le Roi est entré en Boheme ».
%   v. 35     f. 27r    Lalande, « au collège de france le 4 nov. 1797 », a
%                       small slip: an apology after a visit that was not
%                       welcome — « l'indiscretion de ma visite, et
%                       l'improbation de mes hommages » — and a mention of
%                       « mon ami Cousin », of « le systeme du monde de la
%                       place » and of « mon abregé d'astronomie ».
%   v. 37     —         The address leaf of that note: « à la Citoyenne /
%                       Sophie germain / rue Ste croix de la Bretonnerie
%                       n° 18 ». No pencilled folio is legible.
%   v. 38–39  ff. 29r–  Libri, « Florence ce 17. Novembre 1825 », signed
%             30r       « Guillaume Libri »: Arago has carried her a copy of
%                       Riccati; a memoir of his on the theory of heat is at
%                       the Institut under Fourier's eyes; would she
%                       « chatouiller un peu Mr. Fourier » so that Cauchy at
%                       last makes the report he has owed for six months on
%                       the memoir on the theory of numbers.
%   v. 40     f. 30v    The address leaf of that letter, with a postmark.
%   v. 40     f. 31r?   The first page of an unsigned letter in a sixth hand,
%                       opening « Mademoiselle, » — on an agreement reached
%                       between the writer and Germain, « il faut qu'une
%                       chose soit incontestable pour qu'elle s'articule sans
%                       contestation entre vous et moi ». It runs on past
%                       view 40 and is not finished here. No pencilled folio
%                       is legible on this leaf.
%
% Absent from the output: v. 22, 24, 26, 30, 36. Each photographs a blank
% verso through which the facing leaf shows mirrored, and nothing on them is
% right-reading. Their pencilled folios were read where legible — 15 on
% v. 22, 23 on v. 30 — and are used below to name the address leaves that
% face them. The BnF's stamps, the wax seals and the postmarks are not
% transcribed.
%
% Folios. The pencil was read on the rectos at views 21 (14), 23 (16),
% 25 (18), 27 (20), 28 (21), 29 (22), 31 (24), 32 (25), 33 (26), 35 (27),
% 38 (29), 39 (30), and on the blank views 22 (15) and 30 (23). A verso is
% written as its recto's number with v, the leaf being the same. Where no
% pencil could be read — the address leaves at views 21, 25, 27, 37, and the
% leaf at view 40 carrying the unsigned letter — no \folio is given: the
% place in the sequence is plain but it is not written here as though read.
%
% Locating f. 20r. The letter of 18 April 1831 is cited in the literature by
% its folio; this pass read the pencilled 20 at the head of view 27, and the
% four pages are therefore NAF 4073, ff. 20r–21v = Gallica views 27–29.
%
% Notation set once and held to. Fourier writes almost everything in lower
% case, proper names included (« fourier », « prony », « paris »): kept as
% written. Raised endings — « m^de », « m^r », « monsi^eur », « 2^me » — are
% set on the line as « mde », « Mr », « monsieur », « 2me ». « n° » is set
% with the degree sign. Underlined words are set with \emph. A heavy stroke
% that cancels a word is \struck{}; where the cancelled word cannot be read,
% \ill{} stands inside it. On the covers of Fourier's billets the street
% number is written with the 4 and the 1 ligatured and could be read « 4 »
% alone; it is given as \uncertain{41}. Sophie Germain's hand in the letter
% of 1831 is fast and crossed in places by the library's stamp; every word
% this pass could not settle is marked and none is supplied.
\documentclass[11pt,a4paper]{article}
% The preamble shared by every transcription of the mathematicians' manuscripts in Gallica.
%
% It rests on one idea: **uncertainty compounds, it does not resolve.** A
% transcription that smooths over a doubtful word has destroyed the only thing
% it was made for — a reader must be able to tell, at every point, what was on
% the page and what was supplied. The apparatus macros below are therefore the
% critical apparatus, not decoration.
%
% The same commands are recognised by `scripts/render.mjs`, which produces the
% site's reading view, and by `scripts/tei.mjs`. A macro added here must be
% added there too, or rendering fails loudly — which is the wanted behaviour.
%
% Comments here are English, like the rest of the repository. The documents
% this preamble sets are French, and that is deliberate: see the skills.

\ProvidesPackage{germain}[2026/09/15 Transcriptions of mathematicians' manuscripts in Gallica]

\usepackage{amsmath,amssymb,amsthm}
\usepackage{mathrsfs}
% Kept from the parent preamble so the renderer's diagram support stays
% available; these manuscripts draw no commutative diagrams, and nothing here uses it.
\usepackage{tikz-cd}
\usepackage[english,french]{babel}
\usepackage{fontspec}

% --- Greek ------------------------------------------------------------------
% The text face stays fontspec's default, Latin Modern, which has no Greek:
% XeTeX drops every glyph it lacks with a line in the log and nothing on the
% page, and the Greek words the manuscripts quote vanished from the PDFs. So
% the Greek and Greek Extended blocks get a XeTeX character class of their own,
% and entering or leaving a run of it switches the family to CMU Serif — the
% same Computer Modern design, with polytonic Greek — and back. Only the
% family changes: a Greek word in \textit or \textbf keeps its shape and series.
% The transitions to and from every other class carry the tokens that class
% has towards an ordinary letter, so French spacing before « ; » or inside
% « » still holds after a Greek word. No grouping, so no brace or \endgroup
% can come out unbalanced; maths is left alone, since XeTeX inserts nothing
% there. Kept identical in `exercices.sty`.
\makeatletter
\newfontfamily\germainGreekfont{cmun}[
  Extension=.otf, NFSSFamily=germaingreek,
  UprightFont=*rm, ItalicFont=*ti, SlantedFont=*sl,
  BoldFont=*bx, BoldItalicFont=*bi, BoldSlantedFont=*bl]
\newXeTeXintercharclass\germain@greek
\def\germain@greekrange#1#2{%
  \@tempcnta=#1\relax
  \loop
    \XeTeXcharclass\@tempcnta=\germain@greek
    \ifnum\@tempcnta<#2\advance\@tempcnta\@ne\repeat}
\germain@greekrange{"0370}{"03FF}
\germain@greekrange{"1F00}{"1FFF}
\def\germain@greekprev{\rmdefault}
\def\germain@greekin{%
  \edef\germain@greekprev{\f@family}\fontfamily{germaingreek}\selectfont}
\def\germain@greekout{\fontfamily{\germain@greekprev}\selectfont}
\def\germain@greekjoin{%
  \XeTeXinterchartokenstate=\@ne
  \@tempcnta=\z@
  \loop
    \ifnum\@tempcnta=\germain@greek\else
      \XeTeXinterchartoks\germain@greek\@tempcnta=\expandafter{%
        \expandafter\germain@greekout\the\XeTeXinterchartoks\z@\@tempcnta}%
      \XeTeXinterchartoks\@tempcnta\germain@greek=\expandafter{%
        \the\XeTeXinterchartoks\@tempcnta\z@\germain@greekin}%
    \fi
    \ifnum\@tempcnta<4095\advance\@tempcnta\@ne\repeat}
% babel-french sets its own classes, and saves and restores the interchar
% state, each time a language is selected: join again after it does.
\addto\extrasfrench{\germain@greekjoin}
\addto\extrasenglish{\germain@greekjoin}
\AtBeginDocument{\germain@greekjoin}
\makeatother

\usepackage{geometry}
\geometry{a4paper,margin=25mm,marginparwidth=28mm}
\usepackage{marginnote}
\usepackage{xcolor}
\usepackage[normalem]{ulem}
\setlength{\emergencystretch}{3em}
\usepackage{hyperref}
\hypersetup{colorlinks=true,linkcolor=black,urlcolor=black,pdfborder={0 0 0}}

% --- Batch metadata ---------------------------------------------------------
% Taken as they stand by the rendering script, which shows them at the head of
% the reading view. They fix what the file claims to cover. The macro names are
% the parent project's, so that the scripts stay shared; \folder here means the
% volume's slug — `fr-9115` — and \pages counts Gallica views.

\newcommand{\folder}[1]{\gdef\grothFolder{#1}}
\newcommand{\batch}[1]{\gdef\grothBatch{#1}}
\newcommand{\pages}[2]{\gdef\grothFirst{#1}\gdef\grothLast{#2}}
\newcommand{\foldertitle}[1]{\gdef\grothFoldertitle{#1}}

% The holder's own shelfmark, printed as they write it: « Français 9115 ».
\newcommand{\shelfmark}[1]{\gdef\grothShelfmark{#1}}

% The dating, copied verbatim from the holder's catalogue — never inferred
% here. « XIXe siècle » is the BnF's; a pass that dates a leaf from its content
% says so in a \note{}, as its own inference.
\newcommand{\dating}[1]{\gdef\grothDating{#1}}

% The status notice, printed in the title block and repeated as a diagonal
% watermark on every page of the PDF. The manuscripts are in the public
% domain; what this line declares is not a right but a quality — an
% unverified first pass by a machine. The skills write the line; a file
% without it gets no watermark, so leaving it out is visible in the output.
\newcommand{\watermark}[1]{\gdef\grothWatermark{#1}}

\gdef\grothFolder{?}\gdef\grothBatch{}\gdef\grothFirst{?}\gdef\grothLast{?}%
\gdef\grothFoldertitle{}\gdef\grothShelfmark{}\gdef\grothDating{}\gdef\grothWatermark{}

% --- The résumé -------------------------------------------------------------
\newenvironment{resume}
  {\small\setlength{\parskip}{0.45em}}
  {\par\medskip\hrule\medskip}

% --- Keywords ---------------------------------------------------------------
% In English, closing the modernised reading's résumé: the modern vocabulary
% under which to search for what the leaves construct. The manifest extracts
% them to tag the volume — this line is the single source of the tags.
\newcommand{\keywords}[1]{\par\smallskip\noindent
  {\footnotesize\textsc{Keywords} --- \textit{#1}}\par}

% --- The critical apparatus -------------------------------------------------

% The Gallica view number, in the margin. This is the anchor that lets a
% disagreement over a formula be settled by looking at *one* image rather than
% a volume of 750. The site uses it to turn the facsimile as the transcription
% is scrolled. A view is one scanned image: usually one side of a leaf.
\newcommand{\page}[1]{%
  \marginnote{\footnotesize\color{gray!70}\textsf{v.\,#1}}%
  \label{page:#1}\ignorespaces}

% The BnF's pencilled foliation, where the leaf carries it and a pass can read
% it: « 198r ». This is what the literature cites — « ff. 198r–208v » — and
% the one line that turns a citation into a view. Recorded, never inferred:
% a folio a pass cannot read is not written.
\newcommand{\folio}[1]{%
  \marginnote{\footnotesize\color{gray!50}\textsf{f.\,#1}}\ignorespaces}

% The modernised reading's coarser counterpart to \page{}: which views a
% section is drawn from.
\newcommand{\pagerange}[2]{%
  \marginnote{\footnotesize\color{gray!70}\textsf{v.\,#1--#2}}%
  \ignorespaces}

% Illegible: never guessed.
\newcommand{\ill}{\textcolor{red!60!black}{[\,\ldots\,]}}

% A doubtful reading: the word is given, and flagged as uncertain.
\newcommand{\uncertain}[1]{\uline{#1}}

% An editorial addition — a missing letter, an implied word.
\newcommand{\add}[1]{\textcolor{blue!50!black}{[#1]}}

% Struck out by the writer. What was crossed out is part of the document.
\newcommand{\struck}[1]{\sout{#1}}

% The transcriber's note, on the state of the page or the mathematics.
\newcommand{\note}[1]{\par\smallskip{\small\color{gray!60!black}\textit{[#1]}}\par\smallskip}

% A marginal note of hers, distinct from the body of the text.
\newcommand{\marginal}[1]{\par\smallskip\noindent\hspace{1em}%
  {\small\color{gray!60!black}#1}\par\smallskip}

% --- Title ------------------------------------------------------------------
% The title block is French, like the documents themselves.

\AddToHook{shipout/background}{%
  \ifx\grothWatermark\empty\else
    \begin{tikzpicture}[remember picture, overlay]
      \node[rotate=52, scale=3.4, align=center, text=gray!18,
            font=\sffamily\bfseries]
        at (current page.center) {\grothWatermark};
    \end{tikzpicture}%
  \fi}

\AtBeginDocument{%
  \begin{center}
    {\large\bfseries Manuscrits de mathématiciens (Gallica) ---
      \ifx\grothShelfmark\empty\grothFolder\else\grothShelfmark\fi}\\[2pt]
    {\itshape\grothFoldertitle}\\[4pt]
    {\small vues \grothFirst--\grothLast
      \ifx\grothBatch\empty\else\ (lot \grothBatch)\fi}\\[2pt]
    \ifx\grothDating\empty\else
      {\small\color{gray!60!black}Datation du catalogue : \grothDating}\\[2pt]
    \fi
    \ifx\grothWatermark\empty\else
      {\small\bfseries\color{red!45!gray}\grothWatermark}\\[2pt]
    \fi
    {\footnotesize\color{gray!60!black}Transcription établie d'après les images IIIF de Gallica.
     Source gallica.bnf.fr / Bibliothèque nationale de France.\\
     Les lectures incertaines sont soulignées, les illisibles notées [\,\ldots\,],
     les ajouts entre crochets ; \textsf{v.} numérote les vues de Gallica,
     \textsf{f.} la foliotation au crayon de la BnF.}
  \end{center}
  \vspace{1em}}

\endinput


\folder{naf-4073}
\shelfmark{NAF 4073}
\batch{2}
\pages{21}{40}
\dating{XVIIIe siècle}
\watermark{Lecture automatique\\non vérifiée}
\foldertitle{Correspondance de Sophie GERMAIN.}

\begin{document}

\section{Feuillet d'adresse du billet de Fourier « mardi soir »}

\page{21}
Mademoiselle

Mademoiselle sophie germain

rue de braque n° \uncertain{41}

Paris.

\note{ce feuillet ferme le billet transcrit à la vue 19 du cahier précédent ;
aucun chiffre au crayon n'est lisible sur la feuille}

\section{Fourier à Sophie Germain, « jeudi … mai » : une invitation à déjeuner}

\page{21}\folio{14r}
jeudi mai.

\note{un mot rayé entre « jeudi » et « mai » : \ill{}}

mademoiselle,

j'ai l'honneur de proposer à madame votre mere et à vous de dejeuner dimanche
prochain chez moi avec mde de \uncertain{bresslieur} et mde de
\uncertain{vauborel}. je vous prie d'excuser ce que mon projet a de
témeraire et de n'avoir egard qu'à l'intention, excuse commune de tous les
auteurs médiocres. je n'ai assurement aucun moyen de vous recevoir aussi
dignement que je le voudrois, et je sens que je vous demande un sacrifice :
mais j'en serai bien reconnoissant. aurez vous la bonté de présenter mon
invitation à madame votre mere en la priant en votre nom et au mien de
l'agréer.

monsieur et mde de prony seront avec nous. il n'y aura point d'autres
personnes que mon secretaire et moi et peut etre le neveu de mde de
\uncertain{vauborel} s'il est à paris. nous nous réunissons à midi un quart

je vous envoye aussi une gravure dont j'ai quelques exemplaires. c'est le
portrait d'un de nos prédécesseurs, et l'un des plus illustres fondateurs
d'une science que vous aimez et que vous avez augmentée.

agréez mademoiselle avec l'hommage de mon respect l'invitation que je vous
présente et donnez moi ce nouveau témoignage d'une obligeance si gracieuse à
laquelle vous m'avez habitué.

J. fourier.

\page{23}\folio{15v}
Mademoiselle

Mademoiselle Sophie Germain.

rue de braque n° \uncertain{41} au marais

\section{Fourier à Sophie Germain, « mardi matin » : le dîner de décembre et les entretiens analytiques}

\page{23}\folio{16r}
je suis extremement reconnoissant de votre souvenir et de votre bonté. je me
proposois chaque jour depuis mon retour de la campagne d'avoir l'honneur de
vous voir ; j'en ai toujours été détourné par quelque nouvel incident.

\note{le mot « reconnoissant » est traversé de deux traits, l'un sur les
premières lettres, l'autre sur les dernières ; rien d'autre n'a été écrit à
la place}

je me rendrai tres certainement à l'invitation \struck{\ill{}} de madame
votre mere pour mercredi \ill{} décembre, et je désire qu'elle me permette de
lui offrir mon respect dans le courant de cette semaine.

\note{le quantième est un chiffre repassé, suivi de deux lettres en exposant ;
il pourrait être 8, mais la lecture n'a pas été arrêtée}

il me tarde beaucoup mademoiselle de reprendre nos entretiens analytiques.
car vous y mettez beaucoup d'esprit et une sagacité merveilleuse. je vous
demande pardon de cette remarque parce que les éloges vous déplaisent. je
vous promets de tacher de me corriger.

je vous prie mademoiselle d'offrir à madame votre mere \struck{\ill{}}
d'agréer aussi \struck{\ill{}} mes hommages et mes respects.

J. Fourier.

mardi matin.

\page{25}\folio{17v}
Mademoiselle

Mademoiselle sophie germain

Paris.

\section{Fourier à Sophie Germain, « lundi matin » : une note de Prony, le mémoire de Libri}

\page{25}\folio{18r}
j'ai l'honneur de présenter mes respects à mademoiselle germain, en lui
transmettant une note qui m'a été remise par mr de Prony. je ne connois point
l'objet des remarques contenues dans cette note peut etre donneront elles
lieu à quelques éclaircissement que mademoiselle germain désirera faire
parvenir à mr de prony. c'est dans cette vue que je les lui adresse.

je n'ai point trouvé chez moi le mémoire de mr. libri. l'exemplaire que j'ai
parcouru est vraisemblablement celui de la bibliothèque de l'institut. je
vais aujourd'hui à l'académie et je me ferai remettre cet ouvrage que
j'enverrai à mademoiselle germain.

je lui présente de nouveau l'expression de mon attachement et de mon respect.

J. fourier.

lundi matin.

\page{27}
Mademoiselle

Mademoiselle Sophie Germain

rue de braque n° \uncertain{41} au marais

Paris.

\note{feuillet d'adresse du billet ci-dessus ; aucun chiffre au crayon n'est
lisible sur la feuille}

\section{Sophie Germain à Libri, « Paris ce 18 avril 1831 »}

\page{27}\folio{20r}
Paris ce 18 avril 1831

Monsieur, je suis plus \ill{} affligée qu'étonnée de ce que vous méditez
touchant votre situation présente. je vous supposois une inquiétude dans une
disposition d'esprit bien \uncertain{éloignée} de l'amour exclusif des livres
qui aurait pu vous rendre heureux. vous deviez m'avoir écrit deux fois depuis
votre départ, je n'ai pas reçu d'autre lettre que celle écrite lors de votre
départ de chez Mr Maurice. J'ai fait remettre à Made Cauchy la lettre dont
vous m'aviez chargée.

\note{l'adverbe qui manque est masqué par le timbre de la Bibliothèque}

Dans votre lettre vous me mandiez que je recevrois les \uncertain{Jaux} en
même tems que cette même lettre mais vous ne dites pas par quelle voie, je
vous ai répondu en adressant ma lettre à Florence

\note{abréviation d'un mot commençant par J et finissant par une lettre en
exposant, sans doute « journaux » ; elle revient plus bas}

\page{28}\folio{20v}
Je vous priois de me donner de vos nouvelles \struck{\ill{}} j'étois inquiete
de vous et j'avois bien raison, j'en ai fait demander de tout côté. Je vous
mandai \struck{\ill{}} que je n'avois rien reçu mais, avec beaucoup de
\uncertain{sincérité} que ce n'étoit pas cela qui m'occupoit.

Votre dernière lettre m'est arrivée avant hier, j'ai envoyé hier chez Mde
Renaud elle a son adresse elle a remis de suite, à mon domestique qui en a
donné reçu en mon nom, les \uncertain{Jaux}. françois m'a dit qu'elle
paroissoit fort bien portante ce que je m'empresse de vous mander puisque
vous paroissiez inquiet de \struck{\ill{}} santé. cette dame avoit perdu mon
adresse et après avoir lu votre lettre elle a dit n'avoir pas reçu, plus que
moi, les lettres que vous dites lui avoir écrit.

Ma santé est dans un état affreux, une mort prompte seroit pour moi un
soulagement car je souffre des maux inconnus et qui

\page{28}\folio{21r}
ne me laissent pas un seul moment de repos, je voulois lire au moins le 2me
v. de la revue mais je ne puis, je reste enfermée je ne vois pas plus Mr le
Peintre que mes autres amis si ce n'est St amand (qui est toujours fort
occupé de vous) et ma Soeur. Elle dit pourtant que mon état n'est pas
désespéré mais on m'annonce de longues souffrances.

J'ai reçu de Mr \ill{} le Mr qui écrivent un mémoire de vous et \ill{}, et me
demandoit de vos nouvelles je lui ai donné celle de votre départ de genève.
quant à moi je lui ai mandé qu'il m'étoit absolument impossible d'espérer
profiter de la bonne volonté au moins d'ici à longtems.

\note{la phrase est écrite très vite ; le nom du correspondant et un mot
après « de vous et » n'ont pu être arrêtés, et la construction ne se laisse
pas rétablir}

décidément il y a un tort sur tout ce qui tient aux mathématiques votre
préoccupation celle de Cauchy, la mort de Mr fourier, pour \ill{} cet élève
Gallois qui malgré son impertinence annonçoit

\note{deux mots avant « cet élève » n'ont pas été arrêtés : le premier d'une
ou deux lettres, le second commençant par « che »}

\page{29}\folio{21v}
des dispositions heureuses en a tant fait qu'il a été chassé de l'école
normale, il est sans fortune et sa mere en a fort peu. Rentré chez elle il a
continué envers elle cette habitude d'injure dont il vous a donné avant même
un échantillon après votre lecture à l'académie, la pauvre dame a quitté sa
maison laissant de quoi vivre \uncertain{médiocrement} à son fils et s'est
vûe forcée de se placer dame de compagnie pour \uncertain{satisfaire} à
\ill{}. On dit qu'il devient tout à fait fou et je le crains.

\note{le nom de Galois est écrit « Gallois », avec deux l, ici et nulle part
ailleurs dans le volume pour ce que cette lecture en a vu}

Je vous adresse cette lettre à Marseille, indiquant l'avis que vous en
donnerez à Mde Renaud. je vous \uncertain{demande} \ill{} de vos nouvelles,
je vous prie instamment, je désirerois avoir en même tems celle de Made votre
mere. Quoique je n'aie pas l'honneur de la connoître, il m'est impossible de
n'être pas occupée des chagrins qu'elle ne peut manquer d'éprouver. agréez,
Monsieur, l'assurance d'un constant intérêt que vos talens m'ont inspiré.

Sophie Germain

\section{Lagrange à Sophie Germain, « Paris ce 17 germinal (mercredi) »}

\page{29}\folio{22r}
Paris ce 17 germinal (mercredi)

Lagrange présente ses respects à Mademoiselle Germaine. étant de retour de la
campagne où il a passé quelques jours, il se fait un devoir de la prévenir
qu'il sera à ses ordres le 19 et le 20. (vendredi, et samedi) il ne sortira
pas ces jours-là, à moins que Mlle Germaine n'aime mieux qu'il vienne chez
elle, auquel cas il la prie de vouloir bien l'en avertir.

\note{le quantième est un chiffre à barre ; la lecture « 17 » s'accorde avec
le reste du billet, où le 19 est un vendredi et le 20 un samedi, mais cet
accord est une remarque de la présente lecture et non du document. Lagrange
écrit le nom « Germaine »}

\page{31}\folio{23v}
Mademoiselle

Mademoiselle Sophie Germain

Paris

\section{Lagrange, de Berlin, « ce 10 Juillet 1776 »}

\page{31}\folio{24r}
Voici deux lettres que je prends la liberté de vous envoyer et que je vous
prie de remettre à leurs adresses. Saluez de ma part tous nos Amis communs et
donnez moi surtout des nouvelles de notre Ami Denina, pour qui je prends tout
l'intérêt qu'un mérite tel que le sien peut inspirer.

\note{la feuille commence sans suscription : le premier feuillet de la lettre
manque. Elle n'est pas adressée à Sophie Germain}

Je vous prie aussi de dire au Docteur Cigna que je viens enfin de recevoir
son paquet du 10 Novembre et que je lui répondrai incessamment. J'ai demandé
de ses nouvelles à Ms. Richeri, secrétaire de S. E. le Comte Fontana, qui est
son compatriote et même son parent suivant ce qu'il m'a dit, mais il n'a pu
m'en donner ne l'ayant pas vu \struck{\ill{}} avant son départ de Turin.

Conservez moi votre précieuse amitié dont je fais

\page{32}\folio{24v}
le plus grand cas, et à laquelle je tache de répondre par toute la mienne, et
par les vifs sentimens d'estime et de considération avec lesquels j'ai
l'honneur d'être

Monsieur

à Berlin ce 10 Juillet \uncertain{1776}.

Votre très humble et très obéissant serviteur

De la Grange

\note{le troisième chiffre du millésime se lit 7 ou 5 ; Lagrange étant à
Berlin depuis 1766, la seconde lecture est exclue par la suscription
elle-même, mais le chiffre reste douteux}

\section{Lagrange, de Berlin : Denina, les loix du pays, Galilée}

\page{32}\folio{25r}
Je suis bien faché du désagrement que vient d'avoir notre Ami Denina que
j'aime et que j'estime infiniment. j'espere qu'on aura egard à son merite et
à l'honneur qu'il fait à sa patrie ; mais ce que je n'espere gueres pas c'est
qu'il puisse se corriger de ces petites etourderies, qui dans un pais tel que
le votre ne laissent pas de faire beaucoup de mal.

\note{ni suscription ni signature : la feuille est détachée d'une lettre dont
le reste n'est pas dans ce cahier}

Je crois qu'en general un des premiers principes que doit avoir tout homme
sage, c'est de se conformer strictement aux loix du pais dans lequel il vit,
quand meme il y en auroit de deraisonnables. D'ailleurs j'ai toujours observé
qu'en general les ouvrages qui ont attiré le plus de contradictions et de
tracasseries à leurs auteurs, n'etoient pas ceux qui etoient les plus propres
à leur acquerir une reputation solide, temoin l'Encyclopedie et plusieurs
autres ouvrages françois, et même italiens.

Notre grand Galilée \struck{\ill{}} ne doit pas avoir gloire qu'à ses
decouvertes sur le mouvement et sur les satellites de Jupiter. Ses fameux
dialogues auxquels il a du tous ses malheurs, sont les moins bons de tous ses
ouvrages, et on n'en peut plus soutenir la lecture. Sans cela il auroit vecu
plus heureux, et seroit peut etre devenu encore plus grand par

\note{« ne doit pas avoir gloire qu'à » est écrit ainsi ; un mot rayé précède,
et la construction reste boiteuse}

\page{33}\folio{25v}
d'autres decouvertes. Je gage que le nouvel ouvrage de notre Ami, dont il n'a
que du chagrin est bien au dessous de son histoire d'Italie qui
\struck{\ill{}} ne lui a produit que de la satisfaction. Que ne s'exerce-t-il
dans la carrière de l'histoire, ou il a deja eu tant de succès ? c'est la
partie de la litterature que j'estime le plus, et ou il y a le moins de
danger, pourvu qu'on veuille adopter la devise que devroit prendre tout
historien \emph{sine ira et studio}.

J'ai été autrefois, plus que personne, entiché de ces petitesses, et irrité
des persecutions \struck{\ill{}} auxquelles je voyois que souvent les gens de
lettres exposés ; mais je vous assure que j'en suis bien desabusé. L'age, et
peut etre plus encore le climat ou je vis, m'ont donné un sang froid que je
n'avois pas, et qui me fait voir maintenant bien des choses sous un autre
aspect que celui ou j'avois coutume de les voir.

Je joins aux deux jettons une medaille de l'Imperatrice de Russie frappée à
l'occasion du dernier jubilé Academique, et que j'ai recue de l'Academie. on
m'a dit que la tête de l'Imperatrice est très ressemblante.

\section{Lagrange, de Berlin : les estampes du Roi, et les armées en campagne}

\page{33}\folio{26r}
Monsieur

M. Verney qui est parti d'icy a huit ou dix jours avec M. le Marquis de
Rozignan vous remettra d'abord un paquet de ma part contenant quatre estampes
du Roi, dont deux, une grande et une petite, \struck{\ill{}} sont pour vous,
et dont les deux autres sont destinées pour mon Pere à qui je vous prie de
vouloir bien les remettre.

M. Verney a bien voulu se charger aussi d'un autre paquet pour vous contenant
les deux premiers volumes de l'histoire diplomatique universelle de M.
\uncertain{Weguelin}, mais vous ne le recevrez que lorsque l'équipage de M.
le Marquis de Rozignan sera arrivée. je profiterai des occasions qui se
presenteront pour vous envoyer la suite de cet ouvrage à mesure qu'elle
paroitra.

Je vous remercie d'avance de tout ce que vous voulez bien m'envoyer ; mais je
vous avoue que j'ai maintenant quelque

\page{34}\folio{26v}
inquiétude sur l'envoi de M. \uncertain{Rabbi} à cause des grandes armées qui
sont maintenant en campagne, et qu'on ne peut eviter de traverser pour entrer
dans ce pais à moins de faire un grand detour.

J'ai repondu à tous les articles de votre dernière lettre dans celle que j'ai
inserée dans le paquet que M. Verney vous remettra à son arrivée à Turin ;
ainsi je n'y reviendrai plus ici. Je ne vous donnerai pas non plus des
nouvelles de guerre, ne s'étant jusqu'à present encore rien passé, (du moins
que je sache) entre les differentes armées ; mais comme on assure que le Roi
est entré en Boheme \struck{\ill{}} on ne tardera pas à apprendre quelque
chose d'interessant.

Je vous felicite de n'etre que spectateur de cette grande tragedie qui va se
jouer dans nos quartiers.

\note{la lettre s'arrête là, sans signature ; le feuillet suivant manque}

\section{Lalande à Sophie Germain, « au collège de france le 4 nov. 1797 »}

\page{35}\folio{27r}
au collège de france le 4 nov. 1797

il etoit difficile mademoiselle de me faire sentir plus que vous ne l'avez
fait hier l'indiscretion de ma visite, et l'improbation de mes hommages. mais
il m'etoit difficile de le prevoir je ne puis meme encor le comprendre, ou le
concilier avec les talens que mon ami Cousin m'a annoncés. il ne me reste
donc qu'à vous faire des excuses de mon imprudence ; on apprend à tout age et
les leçons d'une personne aussi aimable et aussi spirituelle que vous se
retiennent plus que celles d'\uncertain{autres}.

vous m'avez dit que vous aviez lû le systeme du monde de la place, et que
vous ne vouliez pas lire mon abregé d'astronomie, comme je crois que
\struck{\ill{}} vous \uncertain{auriez} pu entendre \ill{} sans l'autre ;
j'avois cru \struck{\ill{}} \ill{} que le \uncertain{projet} \ill{} formé de
me temoigner l'indignation \ill{} appercevois. c'est l'objet de mes excuses
et de mes regrets.

salut et respect

Lalande

\note{la moitié basse du billet est écrite très menu, les lignes se touchent
et plusieurs mots sont repassés ou rayés : le sens général se laisse suivre,
mais cette lecture n'a pas voulu suppléer les mots qu'elle ne voyait pas}

\page{37}
à la Citoyenne

Sophie germain

rue Ste croix de la Bretonnerie n° 18

\note{feuillet d'adresse du billet ci-dessus, écrit tête-bêche et traversé
d'un grand paraphe ; aucun chiffre au crayon n'est lisible sur la feuille}

\section{Libri à Sophie Germain, « Florence ce 17. Novembre 1825 »}

\page{38}\folio{29r}
Mademoiselle.

Je ne saurais vous exprimer combien j'ai regretté de devoir partir de Paris
sans avoir l'honneur de vous voir : mais des circonstances malheureuses, et
des craintes très-vives que j'avais sur l'état de la santé de ma mère,
m'avaient tellement alarmé que je devins incapable de rien faire, et je ne
sus que partir — à présent mes craintes se sont dissipées, et je m'empresse
de vous écrire pour avoir de vos nouvelles — Monsieur Arago qui a eu la bonté
de passer quelques jours ici, a bien voulu se charger de vous remettre de ma
part, un exemplaire de Riccati que vous paraissiez désirer de posséder :
j'espère qu'il vous aidera à vous rappeler d'une personne qui a pour vous les
sentiments de l'estime la plus profonde — pourrais-je me flatter,
Mademoiselle, de recevoir de tems en tems de vos nouvelles ?

Vous avez tant de bonté pour moi, que j'espère, vous m'accorderez ma demande,
en songeant au plaisir que vous me procurerez — Parlez-moi de vos travaux, de
vos amusements, de la musique, et de tout ce qui vous intéresse, qui doit

\page{39}\folio{29v}
m'intéresser aussi — Je ne peux vous parler à présent de mes travaux,
parceque je n'ai rien publié, et que j'ai été distrait par mille choses —
J'ai envoyé un mémoire à l'Institut sur la théorie de la chaleur, ou je crois
avoir démontré qu'on peut donner à l'Intégrale d'une équation différentielle
linéaire un nombre quelconque de constantes arbitraires — Monsieur Fourier
qui doit avoir à présent le mémoire sous les yeux \uncertain{pourrait} vous
en parler — J'ai aussi déduit les équations du mouvement de la chaleur de
l'hypothèse des vibrations calorifiques — mais vous comprenez bien,
Mademoiselle, que ces choses là ne peuvent pas s'écrire par lettre —

\note{un trait horizontal traverse « pourrait », sans qu'aucun autre mot ait
été écrit à la place}

J'ai aussi fait quelques observations physiques curieuses, que j'enverrai à
Mr. Arago afin de les publier dans les annales.

Voudriez vous avoir la bonté de chatouiller un peu

\page{39}\folio{30r}
Mr. Fourier afin qu'il fasse hâter le rapport que Mr. Cauchy doit faire
depuis six mois sur le mémoire sur la théorie des Nombres que vous avez lu ?
Je vous en serai infiniment obligé.

Je vous prie de presenter mes respects à Mme votre soeur, et à Mr. Legendre,
et d'agréer les sentimens de ma profonde estime, avec lesquels j'ai l'honneur
d'être

Mademoiselle

Florence ce 17. Novembre 1825

P. S. si vous avez la bonté de m'écrire vous pouvez adresser les lettres à
Mon adresse à Florence sous le couvert de mon cousin le Commandeur de
\uncertain{Humbourg} Directeur général des postes de la Toscane à Florence —
(sans les affranchir)

Votre très humble et très obéiss. ser.

Guillaume Libri

\page{40}\folio{30v}
À Mademoiselle

Mademoiselle Sophie Germain

Rue de Braque n° 41 au Marais

à Paris

\section{Lettre non signée à Sophie Germain : « il faut qu'une chose soit incontestable »}

\page{40}
Mademoiselle,

Je suis pénétré de l'intérêt que vous me montrez, Je vous remercie de cœur
d'un sentiment qui me soutient et me console de bien des ennuis. J'ai
remarqué avec reconnaissance le double soin de me dire ce que j'avais besoin
de savoir. Vos idées confirment les miennes \& J'ai remarqué qu'il faut
qu'une chose soit incontestable pour qu'elle s'articule sans contestation
entre vous et moi. Je m'en tiens en Conséquence à ce qui a été dit, sous la
reserve toutefois des \uncertain{menues} dont vous faites les honneurs à ma
sagacité. J'espère à la Vérité les decouvrir plus par mon habitude de marcher
droit que par des calculs dont Je ne tiens pas tous les \ill{}. Je pourrais
donner quelques développemens à \ill{} pensée, mais comme J'espère ne pas
vous être devenu etranger, Je compte sur la facilité avec laquelle vous
achevez ordinairement mes phrases.

Vous avez oublié de me remercier d'avoir été pour vous une occasion de vous
instruire ; sans moi vous n'auriez jamais su qu'il fallait renvoyer les
billets refusés au payement mais à cette petite galanterie près, vous avez
fait tout ce que Je pouvais attendre de votre obligeance \& de votre

\note{la lettre se poursuit au delà de la vue 40 et n'est pas achevée ici ;
elle n'est pas signée sur ce feuillet, et la main n'est celle d'aucun des
autres correspondants de ce cahier}

\end{document}
